% Generated. Do not edit manually. Edit resume_content.json instead.
\renewcommand{\ResumeSummary}{
Ph.D. mechanical engineer combining structural and fluid simulation, nonlinear material modeling, and experimental validation. Developed custom three-dimensional FEM/XFEM and coupled multiphysics solvers, calibrated structural models against measured response, and applied explicit dynamic impact analysis to plastic-component cracking. In product development, used CFD/FEA, parametric automation, and prototype testing to evaluate design alternatives and improve structural and airflow performance.
}

\renewcommand{\ResumeSkills}{
\textbf{Structural Simulation:} FEM/XFEM, Nonlinear material modeling, Explicit dynamic FEA, Modal analysis, J-integral assessment\par
\textbf{Fluid and Multiphysics:} CFD, Navier-Stokes equations, Transient heat transfer, Fluid-structure interaction, Coupled thermal-fluid models\par
\textbf{Numerical Software:} C/C++, Python, CadQuery, Custom solver development, Reduced-order modeling\par
\textbf{Experimental Validation:} Prototype testing, Impact-hammer testing, Structural model calibration\par
}

\renewcommand{\RheemBullets}{
\item \textbf{Explicit Dynamic Impact and Fracture Assessment}
\begin{itemize}[leftmargin=1.3em,topsep=1pt,itemsep=1pt]
\item Performed explicit dynamic FEA impact analysis of plastic components to evaluate impact-induced stresses and potential cracking; calculated the J-integral to assess the driving force for crack propagation.
\end{itemize}
\item \textbf{Multi-Material Product FEA Analysis}
\begin{itemize}[leftmargin=1.3em,topsep=1pt,itemsep=1pt]
\item Performed finite element analysis (FEA) to support product designs using ABS and ASA plastics, aluminum and steel sheet metal, and expanded polypropylene (EPP) at different densities.
\end{itemize}
\item \textbf{Structural FEA and Assembly Optimization of Modular AHU}
\begin{itemize}[leftmargin=1.3em,topsep=1pt,itemsep=1pt]
\item Performed comprehensive finite element analysis (FEA) of the nine-segment modular AHU assembly to evaluate structural integrity, press-fit connections, and load-bearing capacity. Analyzed deformation under fan and motor loads to maintain critical clearances between the rotating fan lip and stationary bellmouth, preventing mechanical interference and supporting reliable prototype operation.
\end{itemize}
\item \textbf{Motor-Mount Structural Optimization}
\begin{itemize}[leftmargin=1.3em,topsep=1pt,itemsep=1pt]
\item Redesigned a motor-mount structure using FEA, reducing material usage by more than 50\% while maintaining stiffness and vibration requirements. Simplified injection-molded components to reduce secondary manufacturing operations and assembly complexity.
\end{itemize}
\item \textbf{Fan-Induced Vibration and Modal Analysis}
\begin{itemize}[leftmargin=1.3em,topsep=1pt,itemsep=1pt]
\item Performed FEA-based modal analysis to identify natural frequencies of product components and assess resonance risk by comparing them with fan rotational frequency (RPM/60) and blade-passing frequency.
\end{itemize}
\item \textbf{Outdoor Fan-Shroud Optimization}
\begin{itemize}[leftmargin=1.3em,topsep=1pt,itemsep=1pt]
\item Used CFD to optimize shroud geometry and improve airflow distribution across the condenser coil. Fabricated and tested a prototype, demonstrating more than 80\% improvement in airflow efficiency (CFM/W) relative to the previous design.
\end{itemize}
\item \textbf{Parametric Fan-Shroud Generation}
\begin{itemize}[leftmargin=1.3em,topsep=1pt,itemsep=1pt]
\item Developed a Python/CadQuery-based tool that converts fan diameter, coil geometry, component locations, and design clearances into 3D shroud geometry and STEP files, reducing repetitive CAD modeling and accelerating design iteration across product configurations.
\end{itemize}
\item \textbf{Refrigerant-Tube Routing Optimization}
\begin{itemize}[leftmargin=1.3em,topsep=1pt,itemsep=1pt]
\item Developing an FEA-based tool to evaluate natural frequencies and optimize compressor-to-outdoor-coil tubing routes to reduce vibration susceptibility. Incorporating experimental modal testing and model updating to calibrate FEA predictions against measured vibration characteristics.
\end{itemize}
}

\renewcommand{\LTTSBullets}{
\item \textbf{Modular GUI for Aerodynamic Data Analysis}
\begin{itemize}[leftmargin=1.3em,topsep=1pt,itemsep=1pt]
\item Developed a modular GUI for aerodynamic data analysis with well-defined input/output interfaces, improving maintainability and reducing computational and memory overhead.
\end{itemize}
\item \textbf{Real-Time C/C++ HVAC Plant Model for SiL and HiL}
\begin{itemize}[leftmargin=1.3em,topsep=1pt,itemsep=1pt]
\item Developed a real-time C/C++ executable HVAC plant model for Software-in-the-Loop (SiL) and Hardware-in-the-Loop (HiL) applications, combining finite element solution of transient heat transfer and velocity-pressure-coupled Navier-Stokes equations with lumped-capacitance heater/cooler models, convective thermal-fluid coupling, and a proportional-derivative (PD) HVAC controller model.
\end{itemize}
\item \textbf{Jenkins CI/CD and Automated Debugging}
\begin{itemize}[leftmargin=1.3em,topsep=1pt,itemsep=1pt]
\item Developed Jenkins CI/CD pipelines with Bash and PowerShell scripting and automated screen-capture capability, streamlining debugging and test workflows for construction and forestry vehicle software.
\end{itemize}
}

\renewcommand{\UTSABullets}{
\item \textbf{XFEM, Nonlinear Materials, and Damage Modeling}
\begin{itemize}[leftmargin=1.3em,topsep=1pt,itemsep=1pt]
\item Developed XFEM methods to represent stress fields around three-dimensional voids and implemented nonlinear material behavior incorporating plasticity, strain-rate effects, thermal expansion, and von Mises and Drucker-Prager yield criteria.
\end{itemize}
\item \textbf{Custom 3D Multiphysics Finite Element Solvers}
\begin{itemize}[leftmargin=1.3em,topsep=1pt,itemsep=1pt]
\item Developed custom three-dimensional finite element solvers from first principles for structural mechanics, transient heat transfer, and fluid dynamics, including coupled structural-thermal habitat models and interior thermal-fluid simulations.
\end{itemize}
\item \textbf{Experimental Structural Digital Twin and Model Updating}
\begin{itemize}[leftmargin=1.3em,topsep=1pt,itemsep=1pt]
\item Created a structural digital twin of an experimental habitat by acquiring impact-hammer and accelerometer data, identifying natural frequencies and mode shapes, using Crypton optical geometry measurements, and calibrating a parameterized beam-element model to match identified modal properties and measured acceleration histories.
\end{itemize}
\item \textbf{Reduced-Order Modeling and High-Order Elements}
\begin{itemize}[leftmargin=1.3em,topsep=1pt,itemsep=1pt]
\item Developed modal-coordinate reduced-order models and applied Guyan static condensation; implemented very high-order finite element formulations for acoustic wave propagation associated with micrometeoroid impacts.
\end{itemize}
\item \textbf{Eulerian CFD and Hurricane Wave Simulation}
\begin{itemize}[leftmargin=1.3em,topsep=1pt,itemsep=1pt]
\item Developed Eulerian-based CFD models to simulate wave generation and fluid motion under extreme hurricane conditions, enabling numerical investigation of hydrodynamic loads on coastal bridge structures.
\end{itemize}
\item \textbf{Coupled Eulerian-Lagrangian FSI Modeling}
\begin{itemize}[leftmargin=1.3em,topsep=1pt,itemsep=1pt]
\item Developed coupled Eulerian-Lagrangian (CEL) simulation models to investigate fluid-structure interaction between hurricane-driven waves and coastal bridge structures, integrating Eulerian fluid modeling with Lagrangian structural mechanics.
\end{itemize}
\item \textbf{NASA Modular-Coupled Virtual Testbed (MCVT)}
\begin{itemize}[leftmargin=1.3em,topsep=1pt,itemsep=1pt]
\item Developed and integrated structural, thermal, and fluid plant models for NASA's Modular-Coupled Virtual Testbed (MCVT), a system-of-systems simulation environment for virtual testing of resilient extraterrestrial habitats.
\end{itemize}
\item \textbf{Efficient Scientific Software and Numerical Algorithms}
\begin{itemize}[leftmargin=1.3em,topsep=1pt,itemsep=1pt]
\item Developed computational tools emphasizing reduced arithmetic and memory complexity to support novel numerical formulations and simulation workflows; contributed to a publication on algorithm development.
\end{itemize}
}

\renewcommand{\OtherExperience}{

}

\renewcommand{\Demonstrations}{
Large-deformation beam simulator: \url{https://www.linkedin.com/feed/update/urn:li:activity:7429372567318331392/}\par
CFD and bridge structural analysis: \url{https://youtu.be/DTQD6CamibM}\par
MCVT structural-mechanical model: \url{https://youtu.be/TAA66P_UXIM}\par
}

\let\ResumeDemonstrations\Demonstrations\renewcommand{\Demonstrations}{}\renewcommand{\OtherExperience}{}
